
\documentclass[11pt,a4paper]{article}

% Encodage
\usepackage[utf8]{inputenc}
\usepackage[T1]{fontenc}
\usepackage{lmodern}

% Mise en page
\usepackage[a4paper,margin=2.5cm]{geometry}

% Tableaux
\usepackage{booktabs}
\usepackage{float}
\usepackage{array}

% Liens
\usepackage[colorlinks=true,
            linkcolor=blue,
            urlcolor=blue]{hyperref}

% Titre
\title{\textbf{Évaluation de l'intermédiation financière
des pensions alimentaires}\\[0.3cm]
\large Données administratives de la CNAF
et faisabilité empirique}

\author{}
\date{\today}

\begin{document}


\maketitle

\section{Évaluation de l'intermédiation financière
des pensions alimentaires : données disponibles
et faisabilité }


\subsection{Présentation de la base ALLSTAT}

\subsubsection{Source et couverture}

ALLSTAT est une base statistique administrative
constituée à partir des fichiers de gestion des
allocataires des caisses d'allocations familiales CAF.
Elle contient des informations relatives à la
composition des foyers, aux ressources prises en
compte dans le calcul des prestations sociales,
aux droits ouverts et aux montants des prestations.

Les fichiers ALLSTAT-FR6 sont référencés dans
le CASD et vont de 2017 à 2025 ie couverture temporelle  adaptée à notre projet, puisqu'elle permet d'envisager l'observation
des allocataires avant et après les deux principales
dates de généralisation de l'IFPA, en mars 2022
et janvier 2023.

L'unité d'observation envisagée est le dossier
d'allocataire à une date de référence donnée.
Cette unité ne correspond toutefois pas
nécessairement à un parent, à un couple séparé
ou à une obligation alimentaire.

\medskip

\noindent
\textbf{Documentation :}

\begin{itemize}
    \item Catalogue CASD :
    \url{https://www.casd.eu/source/fichier-statistique-des-allocataires-des-caf-donnees-definitives/}

    \item Documentation ALLSTAT :
    \url{https://utilcnaf.site.ined.fr/fr/fichiers/allstat}
\end{itemize}

\subsubsection{Population étudiée}

La population d'intérêt serait constituée
des parents séparés ayant au moins un enfant
à charge et concernés par une obligation
de pension alimentaire. Dans ALLSTAT, une première approche consisterait à identifier les foyers
monoparentaux à partir des variables
de situation familiale et de composition
du ménage.

Cependant, l'observation d'un foyer
monoparental ne garantit pas l'existence
d'une pension alimentaire fixée par
un titre exécutoire. Inversement, certains parents séparés
peuvent ne pas être présents dans les
fichiers CAF ou ne pas être identifiables
comme bénéficiaires d'une pension.

L'échantillon effectivement observable
pourrait donc différer de la population
des parents concernés par la réforme.

\subsection{Variables pertinentes pour le projet}

L'examen du dictionnaire des variables
ALLSTAT fait apparaître plusieurs
informations directement utiles
à notre question de recherche.

\subsubsection{Pensions alimentaires}

\begin{table}[H]
\centering
\caption{Variables relatives aux pensions alimentaires}
\small
\begin{tabular}{p{2.4cm}p{10.6cm}}
\toprule
\textbf{Variable} & \textbf{Description} \\
\midrule
\texttt{MTPADU} &
Montant des pensions alimentaires dues. \\

\texttt{MTALIMFO} &
Montant des pensions alimentaires reçues
par le foyer. \\

\texttt{MTALIMME} &
Montant des pensions alimentaires reçues
par la femme. \\

\texttt{MTALIMMO} &
Montant des pensions alimentaires reçues
par l'homme. \\

\texttt{MTALIMVE} &
Montant des pensions alimentaires versées
par le foyer. \\

\texttt{MTPALIFT} &
Montant des pensions alimentaires reçues
par le foyer pour le trimestre. \\

\texttt{MTPALIME} &
Montant trimestriel des pensions alimentaires
reçues par la femme. \\

\texttt{TPA} &
Indicateur de pension alimentaire. \\

\texttt{NBDEBPEN} &
Nombre de personnes débitrices au titre
de l'allocation de soutien familial. \\
\bottomrule
\end{tabular}
\end{table}

La présence de ces variables est
particulièrement intéressante.

Elle indique que les fichiers statistiques
de la CNAF contiennent non seulement
des informations sur les prestations
sociales, mais également sur les
pensions alimentaires dues et reçues.

En particulier, la variable
\texttt{MTPADU} pourrait fournir
une mesure des obligations alimentaires,
tandis que les variables
\texttt{MTALIMFO} et
\texttt{MTPALIFT} pourraient renseigner
les ressources déclarées au titre
des pensions alimentaires.

Toutefois, leur interprétation
nécessite plusieurs vérifications.

Premièrement, il convient de déterminer
si les montants enregistrés correspondent
aux versements effectivement réalisés
ou aux ressources déclarées pour
le calcul des prestations.

Deuxièmement, les périodes de référence
peuvent différer : certaines variables
portent explicitement sur un trimestre,
tandis que d'autres ne précisent
pas leur périodicité dans le dictionnaire.

Enfin, rien ne garantit à ce stade
que les pensions dues et reçues
correspondent à une même obligation
alimentaire.

Il serait donc prématuré de calculer
directement un taux d'impayés
à partir de la différence entre
\texttt{MTPADU} et
\texttt{MTALIMFO}.

\subsubsection{Allocation de soutien familial}

L'allocation de soutien familial (ASF)
constitue un résultat particulièrement
pertinent pour notre étude.

Elle peut notamment intervenir
lorsqu'une pension alimentaire
est absente, insuffisante
ou impayée, selon les conditions
prévues par la réglementation.

\begin{table}[H]
\centering
\caption{Variables relatives à l'ASF}
\small
\begin{tabular}{p{2.4cm}p{10.6cm}}
\toprule
\textbf{Variable} & \textbf{Description} \\
\midrule
\texttt{ASFVERS} &
Allocation de soutien familial versable. \\

\texttt{ASFCOMPL} &
ASF complémentaire versable. \\

\texttt{MTASFVER} &
Montant de l'ASF versable. \\

\texttt{MTASFCPL} &
Montant de l'ASF complémentaire versable. \\

\texttt{EASFV1} &
Indicateur d'ASF versable pour
le premier enfant du dossier. \\
\bottomrule
\end{tabular}
\end{table}

Ces variables permettraient
potentiellement de mesurer
l'évolution du recours à l'ASF
et des montants de droits
associés à cette prestation.

L'interprétation doit néanmoins
tenir compte de la distinction
entre les montants versables
et les sommes effectivement payées.

Par ailleurs, l'ASF a fait
l'objet d'une revalorisation
importante en novembre 2022.

Cette modification intervient
entre les deux étapes de
généralisation de l'IFPA
et constitue donc une réforme
concomitante à prendre en compte.

\subsubsection{Autres prestations sociales}

La base contient également
des informations permettant
d'étudier les conséquences
de la réforme sur les prestations
sous condition de ressources.

\begin{table}[H]
\centering
\caption{Variables relatives aux prestations sociales}
\small
\begin{tabular}{p{2.4cm}p{10.6cm}}
\toprule
\textbf{Variable} & \textbf{Description} \\
\midrule
\texttt{RSAVERS} &
Revenu de solidarité active versable. \\

\texttt{MTRSAVER} &
Montant du RSA versable. \\

\texttt{PPAVERS} &
Prime d'activité versable. \\

\texttt{MTPPAVER} &
Montant de la prime d'activité versable. \\

\texttt{MTPRESVE} &
Montant total des prestations versées. \\

\texttt{MTPFVERS} &
Montant des prestations familiales versées. \\
\bottomrule
\end{tabular}
\end{table}

Ces informations sont utiles
pour examiner les conséquences
indirectes de l'amélioration
éventuelle du recouvrement.

Une augmentation des pensions
alimentaires effectivement perçues
peut modifier les ressources
prises en compte dans le calcul
de certaines prestations.

L'effet sur le revenu disponible
des parents créanciers peut donc
différer de l'augmentation brute
des pensions reçues.

\subsubsection{Revenus et activité professionnelle}

\begin{table}[H]
\centering
\caption{Variables relatives aux revenus et à l'activité}
\small
\begin{tabular}{p{2.4cm}p{10.6cm}}
\toprule
\textbf{Variable} & \textbf{Description} \\
\midrule
\texttt{MTRNIFOY} &
Revenu net imposable du foyer. \\

\texttt{MTREVBRU} &
Revenus bruts du foyer. \\

\texttt{MTREVACT} &
Revenus d'activité du foyer. \\

\texttt{REVACTME} &
Revenus d'activité de la femme. \\

\texttt{REVACTMO} &
Revenus d'activité de l'homme. \\

\texttt{ACTRESPD} &
Activité du responsable de dossier. \\

\texttt{ACTCONJ} &
Activité du conjoint. \\

\texttt{RUC} &
Revenu par unité de consommation. \\

\texttt{NBUC} &
Nombre d'unités de consommation. \\
\bottomrule
\end{tabular}
\end{table}

Ces variables pourraient permettre
d'analyser l'évolution des ressources
et de certaines dimensions
de l'activité professionnelle
des parents créanciers.

Il convient néanmoins de vérifier
les périodes de référence
et les modalités de déclaration
des revenus.

De plus, les données d'un dossier
CAF ne permettent pas nécessairement
d'observer les revenus du parent
débiteur résidant dans un autre foyer.

L'analyse des réponses
comportementales des débiteurs
pourrait donc nécessiter
un appariement avec d'autres
sources administratives.

\subsubsection{Situation familiale et caractéristiques individuelles}

\begin{table}[H]
\centering
\caption{Variables de composition familiale et de suivi}
\small
\begin{tabular}{p{2.4cm}p{10.6cm}}
\toprule
\textbf{Variable} & \textbf{Description} \\
\midrule
\texttt{SITFAM} &
Situation familiale retenue
pour le calcul des droits. \\

\texttt{PRESCONJ} &
Présence d'un conjoint. \\

\texttt{SEXE} &
Sexe du responsable de dossier. \\

\texttt{DTNAIRES} &
Date de naissance du responsable
de dossier. \\

\texttt{NBENLEFA} &
Nombre d'enfants à charge
au sens de la législation familiale. \\

\texttt{NUMCOMDO} &
Code de la commune du dossier. \\

\texttt{IDANORES} &
NIR anonymisé du responsable
de dossier. \\

\texttt{APPAPAN} &
Indicateur d'appartenance
au panel. \\

\texttt{NBMSIFAM} &
Nombre de mois depuis la dernière
modification de la situation familiale. \\
\bottomrule
\end{tabular}
\end{table}

Ces variables permettraient
de définir des sous-populations
et de construire des contrôles
démographiques et géographiques.

Les identifiants anonymisés
pourraient également faciliter
certaines analyses longitudinales,
sous réserve de leur stabilité
et des conditions d'accès.

Toutefois, l'existence
d'un identifiant anonymisé
ne garantit pas qu'il soit
possible d'apparier directement
les dossiers des deux parents
après leur séparation.

\subsubsection{Informations juridiques}

Une variable mérite une attention
particulière :

\begin{itemize}
    \item \texttt{NATJUG} :
    nature de la décision de justice
    assignant un parent défaillant
    au versement d'une créance alimentaire.
\end{itemize}

Cette variable pourrait
potentiellement apporter
une information sur le
fondement juridique de
l'obligation alimentaire.

Cependant, son intitulé
ne permet pas d'affirmer
qu'elle distingue les divorces
judiciaires des autres
formes de séparation.

Or cette distinction
est indispensable pour
la stratégie d'identification
envisagée.

De même, le dictionnaire
ne permet pas d'identifier
clairement la date exacte
de fixation de la pension.

Ces informations devront
donc être recherchées
dans la documentation détaillée
ou dans d'autres fichiers
administratifs.

\subsection{Utilisation de données CNAF dans la recherche}

L'utilisation de données
administratives de la CNAF
pour l'évaluation des politiques
publiques est déjà documentée
dans la littérature appliquée.

Un exemple particulièrement
pertinent est l'évaluation
de la réforme de 2019
de la prime d'activité
réalisée par l'Institut
des politiques publiques (IPP).

Cette étude mobilise
les fichiers ALLSTAT-FR6
pour la période 2017--2020
et construit un panel
longitudinal d'allocataires.

Elle montre qu'il est
possible d'utiliser
ces fichiers pour analyser
l'évolution des prestations
et des ressources des foyers.

Les auteurs soulignent
néanmoins plusieurs
difficultés de suivi,
notamment lorsque les
identifiants administratifs
changent à la suite
de modifications de
la situation familiale
ou de changements de CAF.

Ces limites sont
particulièrement importantes
dans notre cas, puisque
la séparation des parents
constitue précisément
l'événement à partir duquel
l'exposition à la réforme
est déterminée.

\medskip

\noindent
\textbf{Référence :}

Bozio et al. (2023),
\emph{La réforme de 2019 de la prime
d'activité : analyse sur la possibilité
de mesurer les effets sur l'emploi},
Institut des politiques publiques
et France Stratégie.

\medskip

\noindent
\url{https://www.strategie-plan.gouv.fr/files/files/Actualites/2023-10-30%20-%20Rapport%20IPP%20-%20La%20r%C3%A9forme%20de%202019%20de%20la%20prime%20d%E2%80%99activit%C3%A9/ipp_-_reforme_2019_prime_dactivite_-_octobre_2023_v3.pdf}

\subsection{Limites et difficultés d'identification}

L'examen du dictionnaire ALLSTAT
permet d'identifier plusieurs
limites pour notre projet.

\paragraph{Mesure des impayés}

Les variables relatives aux
pensions alimentaires ne
permettent pas encore de
reconstituer avec certitude
les versements mensuels
effectivement réalisés.

Il faudra déterminer
si les montants enregistrés
correspondent à des
ressources déclarées,
à des obligations théoriques
ou à des paiements constatés.

\paragraph{Identification du traitement}

Aucune variable explicitement
identifiée comme indicateur
d'intermédiation financière
effective ou d'éligibilité
à l'IFPA n'apparaît
dans le dictionnaire transmis.

Il reste donc nécessaire
de vérifier si cette
information existe
dans un autre fichier CNAF
ou dans les données ARIPA.

\paragraph{Dates et catégories de séparation}

La stratégie économétrique
envisagée exploite
les différences de calendrier
de la réforme selon
les procédures de séparation.

Elle nécessite donc
d'observer précisément
la date et la nature
du titre fixant
la pension alimentaire.

Ces informations ne
sont pas clairement
identifiées dans
le dictionnaire ALLSTAT.

\paragraph{Construction du groupe de contrôle}

L'analyse suppose
de comparer des parents
soumis à l'intermédiation
automatique à des parents
relevant encore
du régime antérieur.

Il est donc indispensable
de disposer d'informations
sur les obligations
non intermédiées.

Une base limitée
aux seuls dossiers
effectivement pris
en charge par l'ARIPA
ne suffirait pas
à constituer un groupe
de contrôle complet.

\paragraph{Appariement et suivi longitudinal}

Le suivi des familles
après une séparation
peut être compliqué
par la création ou
la modification de
dossiers administratifs.

Par ailleurs,
l'appariement entre
les parents débiteurs
et créanciers
n'est pas garanti
par les identifiants
présents dans ALLSTAT.

\paragraph{Accès aux données}

Les fichiers ALLSTAT
sont des microdonnées
administratives protégées.

Leur utilisation
dans le cadre
d'un projet de recherche
nécessite une procédure
d'accès adaptée.

La possibilité
d'obtenir les variables
nécessaires ainsi que
des appariements
avec les données ARIPA
ou judiciaires reste
à confirmer.

\subsection{Conclusion et prochaines vérifications}

L'examen du dictionnaire
ALLSTAT constitue
une première étape
encourageante pour
la faisabilité empirique
du projet.

La base contient
des informations
sur les pensions
alimentaires dues
et reçues,
les prestations sociales,
les ressources des foyers
et leur composition.

Elle paraît donc
particulièrement pertinente
pour analyser
les conséquences
économiques de la réforme
sur les parents créanciers.

Cependant, elle ne
permet pas encore
de garantir la mesure
des impayés effectifs
ni l'identification
des parents concernés
par les différentes
étapes de la réforme.

Trois vérifications
apparaissent prioritaires :

\begin{enumerate}
    \item Déterminer la définition
    et la période de référence
    des variables
    \texttt{MTPADU},
    \texttt{MTALIMFO},
    \texttt{MTPALIFT}
    et \texttt{TPA}.

    \item Vérifier si les
    données CNAF ou ARIPA
    permettent d'identifier
    la date de fixation
    de la pension,
    le type de procédure
    et le statut IFPA.

    \item Étudier les possibilités
    d'appariement entre
    les dossiers CAF,
    les données ARIPA
    et les données fiscales,
    afin d'observer
    les deux parents
    et leurs ressources.
\end{enumerate}

En définitive, ALLSTAT
offre une source
administrative prometteuse,
mais la mise en œuvre
de la stratégie
d'identification proposée
dépendra principalement
de la disponibilité
d'informations individuelles
sur les obligations
alimentaires,
les paiements effectifs
et les dates d'exposition
à la réforme.

\end{document}
